% =====================================================================
%  R. Nielsen
%  Hr. Professor Brøchners philosophiske Kritik gjennemseet
%  / Professor Brøchner's Philosophical Critique, Examined
%  Kjøbenhavn: Gyldendal (F. Hegel), trykt hos J. H. Schultz, 1868. 52 pp.
%
%  TRANSLATION (English)
%  Translated from transcription.tex (Danish), which is COMPLETE and
%  collated at the image (see its header §4). Mirrors it 1:1: every
%  printed page carries % --- p. N --- and \opage{N} at the corresponding
%  point in the English. Method: ../../../TRANSLATION-PLAYBOOK.md.
%
%  ---- WHAT THIS PAMPHLET IS ----
%  The middle round of the 1868 exchange: Brøchner's _Problemet om Tro og
%  Viden_ -> this reply by Nielsen -> Brøchner's _Et Svar til Professor R.
%  Nielsen_ (texts/brochner/svar-nielsen). Continuous polemic, no heads;
%  closing rule on p. 52. Nielsen quotes Problemet at length.
%
%  ---- TERMINOLOGY ----
%  Deliberately ALIGNED WITH THE OTHER TWO PIECES OF THE EXCHANGE
%  (../../brochner/problemet-tro-viden/translation.tex and
%  ../../brochner/svar-nielsen/translation.tex), so the three can be read
%  side by side and quotations match across them. NB this departs from
%  the older Nielsen translations, which often render Erkjendelse as
%  "knowledge"; here the pamphlet's own distinction Erkjendelse / Viden
%  is kept:
%    Erkjendelse = cognition · Viden = knowledge · Videnskab = science
%    Tro = faith · Villie = will · Bevidsthed = consciousness
%    Forestilling = representation · Personlighed = personality
%    (absolute) Ueensartethed = (absolute) heterogeneity
%    Samvittighedsspørgsmaal = question of conscience
%    Videnskabsspørgsmaal = question of science
%    Miraklet / Underet = the miracle
%  Quotations from Problemet follow the wording of the English
%  Problemet where the Danish is identical.
%
%  ---- CONVENTIONS ----
%  „…“ -> ``…''; \emph{} (letterspacing) carried over one for one;
%  Latin (roman in the print, unmarked in the Danish file) stays Latin,
%  untranslated; ɔ: -> i.e.; Danish work titles stay Danish; footnote
%  marks *) **) with a per-page reset, as in the Danish; the p. 50 note
%  runs over onto p. 51 and is set whole at its call in both files.
%  Printer's errors: the Danish carries them as printed; the English
%  translates the SENSE and logs the divergence in a % note at the site.
%  Brøchner/Bröchner: the print has both; the English writes Brøchner.
% =====================================================================
\documentclass[11pt]{book}
\usepackage[utf8]{inputenc}
\usepackage[T1]{fontenc}
\usepackage{libertinus}
\usepackage{libertinust1math}
\usepackage{textalpha}
\usepackage{graphicx}
\usepackage[english]{babel}
\usepackage[protrusion=true,expansion=true]{microtype}
\usepackage{enumitem}
\usepackage{fancyhdr}
\setlength{\headheight}{28pt}
\addtolength{\topmargin}{-13.5pt}
\usepackage{hyperref}
\hypersetup{
  pdftitle={Professor Brøchner's Philosophical Critique, Examined},
  pdfauthor={R. Nielsen},
  hidelinks
}
\pagestyle{fancy}
\fancyhf{}
\fancyhead[LE,RO]{\thepage}
\fancyhead[RE]{\textit{Professor Brøchner's Philosophical Critique}}
\fancyhead[LO]{\textit{Examined by R. Nielsen}}
\renewcommand{\thefootnote}{\ifcase\value{footnote}\or *)\or **)\or ***)\else
  \arabic{footnote})\fi}
\usepackage{setspace}
\setstretch{1.3}
\usepackage{marginnote}
\newcommand{\opage}[1]{\marginnote{\footnotesize\textit{[#1]}}}
\newcommand{\apage}[1]{\marginnote{\footnotesize\textit{[#1]}}}

\begin{document}

\frontmatter
\begin{titlepage}
\centering
\vspace*{3cm}
{\Large Professor Brøchner's}\\[1em]
{\Huge Philosophical Critique}\\[2em]
{\large \emph{Examined}}\\[1.5em]
{\small by}\\[1.5em]
{\large R. \emph{Nielsen}.}\\[5em]
% [publisher's device: roundel with monogram, not transcribed]
Copenhagen.\\[0.5em]
Published by the Gyldendal Bookshop (F. Hegel).\\[0.3em]
{\footnotesize Printed by J. H. Schultz.}\\[0.5em]
1868.\\[3em]
{\small Translated from the Danish by Hans Halvorson}
\end{titlepage}

\mainmatter

% --- p. 3 ---
\opage{3}\setcounter{footnote}{0}%
In a recently published work, \emph{Problemet om Tro og Viden}, Professor
Brøchner, after having treated and refuted a number of other theories, has
at last, out of statements borrowed from some of my occasional writings and
from another's report, composed for himself a Something that he calls my
``theory''. By subjecting this Something to a detailed philosophical
critique, the Professor has, with perceptible self-esteem, weighty reasons
and many proofs, assured himself that my ``theory'' is good for nothing, and
has recommended his own as the only right one. Since on this occasion I now
come to examine the Professor's philosophical critique, my task will thus be
to find out which components of the theory ascribed to me I can acknowledge
as mine, and how matters stand with the many proofs and weighty reasons.

To understand a critic, one must above all take note of his standpoint. The
standpoint on which Professor Br.\ has placed himself is, naturally, ``the
strictly scientific''. In what he calls the theory's consequences, with
respect to the psychological (pp.\ 178--195), with respect to the ethical
(pp.\ 195--199) and with respect to the religious (pp.\ 199--210), the
critical author has heaped up a number of supposed and real
% --- p. 4 ---
\opage{4}\setcounter{footnote}{0}%
contradictions, which are clearly meant to prove that the theory of faith I
have set up must be rejected when one applies the measure of knowledge. What
has now been achieved by this? Precisely what I wish. Far from undermining
my doctrine of faith, the critique has contributed its share to confirming
just what I inculcate again and again on every occasion: that the doctrine of
faith cannot possibly stand before the tribunal of science, and that it
therefore neither can nor will submit to the verdict of knowledge. To that
extent the Professor, in consequence of a scientific self-deception, has
come to furnish a confirmation where he meant to perform a refutation. But
however much I appreciate the confirmation contained in the critique's
refutation, I must at the same time protest just as emphatically against the
``arbitrariness'' with which the critical author has omitted to take account
of what is peculiar to the theory he means to refute. In the doctrine of
faith I have sketched\footnote{My philosophy of religion has not yet
appeared.} faith is, of course, by no means conceived as an acceptance, supported by
blind authority and obscure feeling, of these or those supernatural facts,
but by inwardly conscious appropriation, personal self-understanding; faith
is, on my conception, not thoughtless but essentially self-thinking, an
original principle of life and for that very reason a distinctive principle
of thought.

It is to faith, not as immediate acceptance but as a principle of thought,
that I have ascribed a \emph{critique} of its own, in that I have
distinguished sharply between the \emph{critique of faith} and the
\emph{critique of knowledge}. What, now, has Professor Brøchner's critique
brought to light on this point? Two formal objections, of which the first
cancels itself because it is superfluous, and the second falls away because
it is unfounded. The first (pp.\ 190--91) is to the effect that the
critique of knowledge cannot acknowledge the critique of faith, and about
this there is
% --- p. 5 ---
\opage{5}\setcounter{footnote}{0}%
no dispute at all; the second is to the effect that the critique of faith
conflicts with the essence of faith, but it is without foundation. In proof
that it is without foundation, we let the Professor speak for himself
(pp.\ 192--93): ``When e.g.\ the account of Paradise `with all
character-traits contained therein' is said to have essential and true
reality, then for knowledge an `inner' reality, a `higher reality', in which
all the character-traits contained in that account are preserved, is only a
fantasy-reality, and faith will not be content with a sublimated reality
that lacks the facticity on which faith as faith must lay decisive weight.
Faith demands that the factual congrue with the ideal, with the faith-ideal,
that reality in this sense become a higher reality, but it protests in the
most determinate way against that doubling of reality which deprives faith's
object of the real significance of the historical. Were one thus to let
Christ's birth, death and resurrection take place, not in the outer
historical reality in which human life has its stage, but only in that
higher reality to which they, as falling under the concept of the miracle,
must according to the theory belong, then faith would have to see herein an
annihilation of the whole foundation of Christianity. And science would only
be able to grant it right in this, just as science on its own behalf in
general, while it protests against being withdrawn from the true reality of
ideality, must just as determinately protest against a doubling that
falsifies the reality of ideality by setting in its place a bastard-reality,
born of fantasy, and one which moreover must become arbitrary and
inconsistent. Every miracle would have to acquire that higher reality for
its element, but of the miracles that have an external aspect one may only
trace those that can be reduced to visions, to `expressions' of the
believer's inner relation to the divine, back
% --- p. 6 ---
\opage{6}\setcounter{footnote}{0}%
to it; at least one may only do so unreservedly with respect to them. Thus
e.g.\ the miracle of the burning bush can become an inner reality, but what
now of such miracles as e.g.\ the crossing of the Red Sea, or the springing
forth of the water from the rock at Moses's blow? These one may not deprive
of the actuality of outer reality, and why not? Perhaps on a ground of
faith? No, on an inconsistent, from faith's standpoint arbitrary, rational
regard for the historical facts, which indeed speak for the Jews' having
crossed the Red Sea in the sense of outer reality, and for the thirsting
wanderers in the desert's not having had to be content with an `inner',
`subjective' water, or a `to-come' water, but having been saved from
perishing by a water actual in the sense of outer reality, which faith still
% Danish prints "man reddedes" where Problemet (p. 193) has "men reddedes"; translated in the sense of Problemet's "men" (but)
seeks in an actually existing spring''\footnote{As a note is added:
``Cf.\ 4 Moses XX, 11: and they drank thereof, and their cattle.''}.

This piece of critique is remarkable in two respects.

In the first place, it can be seen from it that the theory the Professor
attacks is not my theory. The concept of faith he takes as his basis departs
from mine to such a degree that I will not waste a word on illuminating the
departure. The conception of the miracle that he ascribes to me
``congrues'' with mine as the square with the circle; anyone who will merely
cast a glance into my writings can convince himself of that. Since I
conceive the concept of the miracle in the most determinate way not as a
concept of knowledge but as a concept of faith, I have no need or incentive
whatever to have objective reality explained away where it really forms an
indispensable moment of the phenomenon of faith. It is in truth not I who
have tried to let either the Israelites or their cattle quench
% --- p. 7 ---
\opage{7}\setcounter{footnote}{0}%
their thirst in subjective water: Professor Brøchner's ``subjective water''
springs solely from his own ``spring''.

In the second place, we learn from the above statement how immature,
thoughtless and helpless faith is after all when it is forsaken by lofty
science. Poor science-less faith, then, has not even so much judgment that
it can distinguish, in a biblical account, between outer and inner reality,
between spirit and letter; and indeed where should it get judgment from,
when all judgment comes into being in science and belongs to science? From
this one can well explain to oneself that faith, ignorant of the critique of
knowledge, must, as far as the account of Paradise is concerned, either take
everything literally or reject everything. To distinguish between lower and
higher reality, between material reality and reality of spirit, reality for
the spirit, requires, after all, sound understanding; but where should faith
left to itself get sound understanding from, when all understanding is a
fruit of science and belongs to science? Faith as faith may, in the eyes of
objective knowledge, do what it will: it is lost. If it attempts by a
critique to separate out what does not concern its essence, then it has
already, after all, laid hands on ``reason'', has in the same act begun to
``rationalise'' itself and prepared its own downfall in science. If it
spurns all critique and grasps blindly at what appeals to its ``obscure,
selfish drive'', then it betrays of itself its innate thoughtlessness;
objective knowledge, to put it briefly, will hear nothing of thoughts of
faith: faith is stupid and shall be stupid, and there's an end of it.

If we ask what entitles philosophical critique to assert so unconditionally
the sole right of knowledge in the world of spirit, then to find the right
answer we must naturally go back to ``the foundation''. The author's
% --- p. 8 ---
\opage{8}\setcounter{footnote}{0}%
only light is the light of science; for faith is only a mystery, and the
mystery is outer darkness. But when science is allowed without further ado
to reject every measure other than its own; when all negotiations about
faith and knowledge are conducted from the outset exclusively from the
standpoint of knowledge, when everything is decided in knowledge and by
knowledge: what wonder, then, that ``the Problem of Faith and Knowledge''
comes out in favour of knowledge! It does indeed look as if science, through
the critique, were to make good its right to dissolve faith, positive faith
in the miracle of revelation; but that is a semblance. Science does not
prove its right; it only postulates it; for every proof is a proof of
knowledge, thus a proof that presupposes precisely what should first of all
be proved, namely that the proof of knowledge has validity in the domain of
faith: the fundamental proof is lacking, and, since the fundamental proof is
lacking, the whole demonstration moves alternately between petitio principii
and circulus in probando.

Right at the beginning (p.\ 130) there is put forward, to my great surprise,
the assertion that, in order to maintain ``the opposition between existence
(life) and knowledge (theory)'' and between ``subjective and objective'', I
take the concept of existence in ``an arbitrary double meaning'' and build
``a sophistical inference'' upon it. --- ``\emph{Existence}'', it is said,
``is taken, \emph{first}, in order to be asserted as a universal
presupposition and necessary condition for every concrete determination,
thus also for knowledge, in the meaning of \emph{individual being in
general, and then}, in order that through the concept of existence the
practical may be claimed a priority and a preponderance, in the meaning of
\emph{volitional being} (`existence is practical'). Without this subreption
the argumentation, which is supposed to lead to the conclusion that it is
possible to place existence highest and yet retain theory, but impossible to
place theory highest
% --- p. 9 ---
\opage{9}\setcounter{footnote}{0}%
and yet retain existence, and further to establish the superiority of faith
over knowledge, could not advance at all. And in the proof of this
proposition, the first premise---a consciousness can exist without being
theoretical, but cannot be theoretical without existing---also contains an
untruth. All being in the form of individuality is existence; there can
therefore be an existence that is not consciousness, not theoretical. But
when \emph{consciousness} becomes the subject of existence, then
\emph{existence includes theory}, for \emph{consciousness is only existing
as being theoretical}''. Tot verba\ldots, we shall now see what lies in
these verba!

That I really have used, and still intend to use, the word existence not
only of the lower, unconscious, but also of the higher, self-conscious mode
of being ``in the meaning of individual being in general'', I grant without
reservation; but how, in a use of a generally familiar word in keeping with
common usage, there can lurk either ``ambiguity'' or ``unclarity'' or
``subreption'', I am unable to see. Wherever the opposition between
\emph{existence} and \emph{theory} is treated in my writings, the word
\emph{theory} is expressly taken as synonymous with the word \emph{science,
scientific objective knowledge}. When I say that a consciousness can exist
without being theoretical, but cannot be theoretical without existing, it in
truth takes a peculiar kind of acuteness to misunderstand the meaning; for
the meaning is that a consciousness can exist without being scientific, but
cannot be scientific without existing. But when this is incontestably and
without any circumlocution the meaning, the clear, obvious meaning, are not
logical spectacles of unusual grinding required to discover in so simple a
% --- p. 10 ---
\opage{10}\setcounter{footnote}{0}%
proposition ``unclarity'', ``arbitrariness'', ``subreption'' and
``sophistry''?

But, it will perhaps be objected, the unclarity is proved after all, and
that already on the first pages (pp.\ 131 ff.), and proved with respect to
the conception both of ``the subjective'' and of ``the objective''. Before
we go further, then, we will pause a little at these categories and consider
each of them more closely by itself.

It is objected that ``subjectivity'' and ``the subjective'' are taken in a
``\emph{double meaning}, so that it now comes to express the accidentally
individual and the egoistic, which falls outside the objectively valid, and
now the universal form of subjectivity''. It is overlooked, says the critical
author, that the form of subjectivity also contains what is objectively
valid, that the subject takes it up as a norm in itself ``without the content
being changed, that subjective action can have objective principles, that
there is an objectively true and valid element in reason-determined
subjectivity. The standpoint of subjectivity, of life, comes in that'' (my!)
``conception to coincide with the standpoint of egoism''. Into this account
the critique has put something that does not concern me. In my conception of
the essence of subjectivity the emphasis by no means falls on the opposition
between the individual and the universal, for every existing subjectivity is
as such individual-universal, but on the opposition between lower and higher
self-regard: the lower self-regard is \emph{egoistic}, the higher is
\emph{religious}. At the stage of egoism, individuality comes into manifold
conflict with the objective powers of existence; the subjectivity moved by
the principle of faith, on the other hand, will know itself reconciled with
the divine laws of nature as well as of spirit, in that it does not give up
but vigorously asserts its individual essence.

% --- p. 11 ---
\opage{11}\setcounter{footnote}{0}%
On my way of seeing it, the objectively universal is so far from being able
to devour the subjectively individual that self-regard, one's own personal
self-regard, on the contrary becomes all the stronger, deeper, more inward,
the higher subjectivity stands in the kingdom of spirit. Should it really be
the case, as Professor Brøchner asserts, that ``the subject takes up what
is objectively valid''---what is objectively valid, be it noted, in
\emph{knowledge and science}---``as a norm in itself, without the content
being changed, that subjective action can have objective''---objective, be
it noted, in the sense of knowledge and science---``principles'', that
otherwise ``the standpoint of subjectivity and of life'' must ``coincide
with the standpoint of egoism'': then, yes then, it follows incontestably
that positive religion must be abolished. Should religious self-regard
really be egoism, faith in a personal God egoism, faith in the immortality
of the soul egoism, faith in revelation altogether egoism, then it certainly
cannot be denied that the Gospel, in spite of self-denial, patience and
love, is originally --- through ``obscure instinct'' and ``selfish drive''
the offspring of concealed egoism, and the hero of Gethsemane the greatest
egoist who has lived on earth.

As proof of my unclear and erroneous conception of the objective, the
following is adduced (pp.\ 132--33): ``In the same way as the concept of
% Danish letterspaces only "jectivitetens" of the line-divided "Ob-jectivitetens"; the English emphasises the whole word
subjectivity, the concept of \emph{objectivity} also receives a \emph{double
meaning}. Where the main stages of subjectivity are
developed''---reference is made here to my Philosophical Propaedeutics,
% "det objective Ydre": rendered "the objective external", as in the English Problemet (corrected there 2026-09-22 from "objective evil")
\S~15---``the objective now means the objective external, or the merely
abstract universal, and now the universal and universally valid that
encompasses the totality of the subjective. And as the abstract opposition
is definitively maintained, it is arbitrarily ignored that at the different
stages the subject's determinateness
% --- p. 12 ---
\opage{12}\setcounter{footnote}{0}%
---which especially appears clearly at the stage designated as that of
infinite subjectivity---is precisely by virtue of what has been determined
as the objective principle of existence, and that it expresses this. It is
overlooked that subjective truth is precisely the appropriated and, in the
appropriation, essentially identical objective truth''.

That the philosophical critic has here, in a rather uncritical manner,
mixed together two quite different kinds of the universally valid, the
subjective and the objective, can easily be shown. It has, strangely enough,
% Danish prints "Opmærhsomhed" (for Opmærksomhed); translated as attention
escaped his attention that what is said at the passage cited from the
Propaedeutics about ``infinite subjectivity'' aims not at proving its truth
but at proving its untruth. It reads: ``In vigorous one-sidedness a choice
is made between two mutually opposed ways out. Partly the objective, in all
its forms, is so reduced to a mere means that the enjoying individuality can
make itself master of things (hedonism, Aristippus); partly the objective, as
the universally rational, is so breathed into consciousness that everything
finite is transformed into something indifferent; purified and hardened in
infinite resignation, the acting subjectivity will determine itself in being
determined by fate (Stoicism; Zeno, Epictetus)''. That what is aimed at here
is a self-deception called forth by the objective is clear from the context,
and is, to superfluity, expressly said in what immediately follows, where the
characterisation of ``\emph{absolute subjectivity}'' begins with the
following words: ``Formal self-liberation ends in the negative;
subjectivity''---the subjectivity that would determine itself by virtue of
the objective principle of existence---``seeks life, but finds death''. It
is more than a distortion of my words, it is a thoroughgoing
misunderstanding of the whole matter that is
% --- p. 13 ---
\opage{13}\setcounter{footnote}{0}%
at issue here. It is overlooked, according to Professor Brøchner's
assurance, that subjective truth---at the stage of hedonism and of
Stoicism---``is precisely the appropriated and, in the appropriation,
essentially identical objective truth''. But now, as is well known, the
Stoic principle stands in decisive conflict with the hedonic; the two
principles, then, cannot possibly be equally objective, unless the objective
validity of principles is to be recognised precisely by their standing in
subjective conflict with one another. The Stoic, certainly, holds himself
convinced that the Stoic principle is the only true, universally valid one,
that every rational human being ought to live as a Stoic; but the hedonist
is, after all, for his part and as a hedonist, just as convinced that the
% „med sig lige objective Sandhed“ (Problemet: „med sig væsentlig lige objective“): rendered
% literally, "equally objective with itself", because Nielsen's retort below turns on „lige objective“
principle of enjoyment expresses ``the appropriated truth, equally objective
with itself in the appropriation'', that the human being who does
not understand how to enjoy does not understand how to live either. If the
hedonist expresses ``the truth, equally objective with itself in the
appropriation'', then the Stoic expresses the untruth that is, in the
appropriation, ``equally objective'' with itself, and vice versa. Who is now
to judge between the contending parties, and from what standpoint is the
judgment to be passed? ``From the standpoint of theory, of science, of
objective knowledge,'' answers the critical author. But if the philosophical
error of the philosophical critique can strike the eye anywhere, it must
surely be here. A Stoic can, as is well known, familiarise himself with the
hedonic theory just as securely as a hedonist; conversely, a hedonist can be
just as much at home theoretically in the Stoic theory as a Stoic, and still
the Stoic does not become a hedonist, nor the hedonist a Stoic. Can there be
a more telling testimony against Professor Brøchner's ``appropriated truth,
equally objective with itself in the appropriation?''

% --- p. 14 ---
\opage{14}\setcounter{footnote}{0}%
How different theories of this nature are from the truly objective theories
of science can be seen from examples in plenty. Think of the Pythagorean
theorem: how should it be possible to familiarise oneself with the proofs of
this theorem and in spite of that maintain a mathematical ``standpoint'' that
stood in open conflict with the Pythagorean? Think of the Copernican system:
how should it be possible to familiarise oneself with the proofs by which
modern science has established the objective validity of this system, and
still maintain an astronomical standpoint that stood in open conflict with
the Copernican? Think of the doctrine of colours in optics, of the doctrine
of the circulation of the blood in physiology; in short, take whatever
knowledge-theories you will, it holds of them all that clear insight into
the objectively valid theory coincides with acceptance of the theory. But
clear insight into a subjectively valid theory by no means coincides, as
Professor Brøchner has so completely overlooked, with acceptance of the
theory. And why not? Because the validity of the subjective theory rests
essentially on subjective grounds, and subjective grounds on the individual
subjectivity.

That the Stoic cannot convince the hedonist, nor the hedonist the Stoic,
however much they work at the theories with each other, is a plain
consequence of the practical, existential opposition of their standpoints.
What in all the world can it help that the Stoic theoretically points out
the supposedly objective contradictions in which the theory of enjoyment can
entangle the one who enjoys, so long as the hedonic subjectivity continues to
posit enjoyment as the highest thing in life? That the contradictions cannot
be resolved objectively in an objective theory teaches him precisely that it
is not the theoretical but the practical on which the
% --- p. 15 ---
\opage{15}\setcounter{footnote}{0}%
emphasis falls, that the practical is essentially subjective, and subjective
practice an art of enjoyment. The art of enjoyment, certainly, cannot resolve
all contradictions either, but what does that prove? ``That the principle of
enjoyment is itself a contradiction!'' By no means. It is precisely the
hedonist's irrefutable maxim that the enjoyment of life is itself the
resolution of contradiction, that the resolution of contradiction is
enjoyment. No, it proves only that life presents contradictions that no one,
the Stoic as little as the hedonist, is able to resolve. When the
contradictions of life become too strong for the Stoic, what does he do then?
He takes his own life. But the same, the very same practical way out stands
open, after all, to the hedonist too. Theory as theory cannot possibly bring
about a change in the subjective standpoint, since the weight of the grounds
that are to decide the preference of the one theory over the other depends
precisely on the practical maxim of subjectivity.

If there really could be an objective insight, independent of every
individual and subjective fundamental disposition, into the original
self-validity of human subjectivity, then it would above all have to be a
matter of concern to science to bring the question of the immortality of the
soul to a decisive answer; for that the question of the immortality of the
soul is inseparable from the question of the individually universal reality
of subjectivity, surely no one will deny. But that science, when it comes to
the point, can no more prove the impossibility of the immortality of the soul
than it is able to prove either its actuality or its necessity, is a settled
matter. He who asserts that the impossibility can be decided by virtue of
``the objective principle of existence'' is just as surely caught in a
self-deception as he who thinks that his conviction of the necessity is owed
to his knowledge. He who denies immortality denies it manifestly by virtue of
% --- p. 16 ---
\opage{16}\setcounter{footnote}{0}%
a subjective, individual fundamental disposition; his pretended grounds of
science are not objectively valid grounds but improbabilities, to which the
subjectively bribed reflection imparts a semblance of objective proofs.
% [p. 16 markers and the run-on sentence are in the previous fragment]
It is not the same with a judgment of the self's reality as with a
judgment of the reality of the rainbow, or of the cloud, or of the
mineral; the reality of the rainbow, the cloud, the mineral is of an
objective nature and can be judged objectively according to the
universally valid laws of knowledge; but the self's reality is of a
subjective nature and must be determined subjectively. The reality of the
objective forces itself upon thought with necessity; the reality of the
subjective announces itself only on the condition of freedom. To want to
derive the reality that ought to be ascribed to one's own individual,
subjective being from the reality that, according to scientific
deliberation, must be ascribed to subjectivity in objective generality, is
a backward and self-contradictory procedure, since the judgment about
subjectivity in general depends precisely on the judgment about one's
own. The validity that a man feels he \emph{must} ascribe to his own
subjectivity rests in the last instance on the validity he \emph{wills} to
ascribe to it; in this sense it holds that the determination of reality is
essentially a determination of will.

The cognition that the will as will must necessarily presuppose is not,
as Prof.\ Brøchner with forced superiority would make out, a theoretical
cognition, an objective cognition of ``the objective principle of existence'';
no, it is just the very opposite, it is precisely a practical cognition, a
subjective cognition of the subjective principle of one's own existence.
It is, in other words, an existential self-cognition; but an existential
self-cognition can be called theoretical only by an arbitrary departure
from all usage of language.

Whether the existential self-cognition stands, for the rest, on a lower or
% --- p. 17 ---
\opage{17}\setcounter{footnote}{0}%
on a higher stage, this much is at any rate certain, that its further
development by no means signifies its transition into objective knowledge;
on the contrary! the purer, clearer and deeper the practical
self-cognition is, the more distinctly it separates itself from the
theoretical. At its first, natural stage self-cognition is immediately
sensuous and the corresponding will egoistic; at the spiritual stage of
freedom, and so at the highest stage of subjectivity, self-cognition is
ideal, and the corresponding fundamental will religious: faith as personal
assurance is no theory, it is an inward, original deed of the spirit, an
act of the unity of cognition, of self-cognition, of practical cognition,
with the fundamental will.

How matters stand with egoism and the cognition of the universal can now
easily be made clear; for it is by no means the case, as Prof.\ Brøchner
supposes, that the cognition of the objective universal should be able to
free subjectivity from its egoism. Insofar as the fundamentally human is
originally in every human being, each single human subjectivity must
undeniably be able to express what is universally valid for human
existence in proportion as cognition clears itself. But that this
universally valid is and remains subjective none the less, its cognition
a decisively subjective, practical cognition, is incontrovertible.
Practical cognition is a cognition of will, a cognition which, in
conditioning the will, is conditioned by the will, a personal
self-cognition. So long as there is something impure, drive-bound and
egoistic in the will, so long is there also something impure, drive-bound
and egoistic in the cognition, and conversely: the movement goes round in
a psychological circle. Stoic egoism is another than the hedonic; but the
Stoic is nevertheless just as essentially an egoist as the hedonist. In
the objective
% --- p. 18 ---
\opage{18}\setcounter{footnote}{0}%
universal of knowledge there is indeed abstraction from egoism, insofar
as there is abstraction altogether from one's own I; but this abstraction
is no personal liberation, no actual redemption: the actual redemption is
only in faith. All objective knowledge is conditioned by the personal
concerns of one's own existing subjectivity being brought into
self-forgetfulness; it is therefore impossible that any objective truth
whatever should be able to raise the knowing subjectivity to spiritually
free self-sufficiency. Where objective knowledge begins, self-regard ends;
where self-regard begins, objective knowledge ends. With returning
self-regard we have the selfish again, and with it the psychological
circle. Objective cognition can be pure, and self-cognition still impure:
deliverance is only in faith. One of two, that is faith's dilemma: either
the unfree, selfish subjectivity must remain eternally in egoistic,
drive-bound self-cognition, egoistic, drive-bound will, or else, on the
condition of higher personal aid, let itself be redeemed by the
subjectivity that is originally pure, freed from drive-bound egoism, and
just for that reason freeing in spirit and in truth.

It is a dangerous thing to confound the subjective with the objective?
The believer sees the universally valid in revelation; the unbeliever
denies revelation by virtue of the universally valid. If now both of
them, the unbeliever as well as the believer, in selfish, passionate zeal
go and confuse the subjective universally valid with the objective, what
have we then? Why, then we have fanaticism. The fanaticism of faith is
well enough known; its mark is that the self-sick subjectivity, in
passionate blindness, would transform subjective truth into an objective
power, a law of existence, before which all other subjectivities are
unconditionally to
% --- p. 19 ---
\opage{19}\setcounter{footnote}{0}%
bow. The fanaticism of knowledge is less well known; but not for that
reason less characteristic. It naturally does not have its source in
science, but in a selfish distortion of science. Instead of taking the
cognition of the personal for what it is, a subjective, practical
cognition, the knower makes self-cognition without further ado
objectively theoretical. It is a great psychological self-deception. The
knower imagines that the subjectivity from which he proceeds is an
objectively universally valid being, an infallible metaphysical
consciousness, a true normal subjectivity, while in actuality it is
nevertheless only an abstraction from the accidental state of his own
individual subjectivity. Far from being freed from ``egoism,'' the knower
is on the contrary infatuated by egoism through this self-deception; for
to make the accidental bent of one's own subjectivity into the objective
standard for all other subjectivities is egoism. What the normal
subjectivity, ``by virtue of the objective principle of existence,'' has a
need for, that every human subjectivity is now essentially to have a need
for, namely to become metaphysical; what the normal subjectivity feels no
need for, a personal God, a comfort in prayer, an eternal life, etc.,
such dreamings and subjective old wives' fancies no true---``objectively
true''---human subjectivity may henceforth feel a need for. To feel a
need for such things is, after all, to deviate from the normal; to
deviate from the normal is to deviate from the universal; to deviate from
the universal is to deviate from the truth; the standard is objective.
With its objective standard the normal subjectivity now needs only---in
all impersonality---to have a little fit of personal zeal for truth, and
the smouldering fanaticism of knowledge will betray itself by an attack
on faith.

From what has been set out here the difference between
% --- p. 20 ---
\opage{20}\setcounter{footnote}{0}%
my well-known way of seeing and the something that Prof.\ Brøchner has
criticized under the name of my ``theory'' must surely be conspicuous. By
holding fast to this difference it then becomes a fairly easy matter to
pass a well-grounded judgment on the scientific worth that can be
ascribed, on the whole, to the Professor's philosophical critique.

The critique has taken its point of departure in psychology. Under the
heading ``\emph{The Psychological Opposition between Will and
Knowledge}'' there is presented a section reflected in fog-veiled
general propositions (pp.\ 133--160), whose aim is to show how
crookedly, arbitrarily and inconsistently I have apprehended the will's
relation to knowledge in my insistence on the essential difference
between practical and theoretical cognition. In the proposition I
uttered: ``what is most subjective of all in subjectivity is the self;
but what is selfish in the self is self-determining, is will,'' the
fundamental error is supposed to lurk.
% "det Selviske i Selvet": Nielsen's play on Selv/selvisk (the self-ish in the self); "selfish" kept for the play
To correct that blunder Prof.\ Brøchner will then set out: ``\emph{from
essence-universality's} determination, which expresses itself in the
being-for-itself in the form of liberated universality that is
self-consciousness's,'' and apprehend ``the essence's fundamental
determination as energy, more closely as selfhood-energy.'' If obscurity
and profundity are one and the same, then the corrective applied here is
undeniably very profound; if foggy generality is a proof of objective
knowledge, it shall be willingly conceded that this psychological
corrective is couched in objective knowledge. But if the task is a
development of the matter, if we, without fear of offending the universal
demands of strict knowledge, may actually venture upon making clear the
peculiar nature of peculiar relations, then I cannot see otherwise than
that the corrective applied with such grand emphasis comes to just about
nothing.

% --- p. 21 ---
\opage{21}\setcounter{footnote}{0}%
What is set up here as a corrective against me is, after all, literally
---even down to the watchwords ``selfhood and selfhood-energy''---my own
view. If anything was to be accomplished here, then, a satisfactory proof
would have to be brought that my view was in conflict, not with Prof.\
Brøchner's peculiar view, for on this point he has no peculiar view, but
in conflict with itself. Now, has such a proof really been brought? We
shall examine it. When I say that what is most subjective of all in
subjectivity is the self, I have by that very token stated a difference
between subjectivity and the self; when I say that what is selfish in the
self is self-determining, is will, I have by that very token stated a
difference between the self and what is selfish in the self; when I say
that what is selfish in the self must, in order to be will in the proper
sense, be self-determining, I have by that very token admitted a
heightening of self-determining, a psychological ascent from the
selfhood-energy's first obscure, involuntary manifestation to the clear,
conscious concentration that expressly marks the will as the act of
selfhood positing itself, asserting itself, resting on itself, i.e.\ as
choosing, resolving, accomplishing will.

If the general could make the matter clear by means of the general in
general, then one would certainly have to grant the critical author
that he acquitted himself well. For there is surely no shortage of
varying assurances about the will in the universal, about the universal
in the will, about the will's unity with the universal in general. ``The
universality that characterizes cognition'' is ``a universality that
strives toward concretion'' etc.\ (p.\ 135). ``Conversely'' (p.\ 136)
``the determination of singularity in the will is such a one as with
necessity unfolds the universal in itself,
% --- p. 22 ---
\opage{22}\setcounter{footnote}{0}%
because the will is only will as the self-conscious, i.e.\ according to
its concept universal and general, essence's will. The will must unfold
itself from the will determined by the single, positing the single, to
the will determined by the universally rational that expresses the
essence, realizing the universal.'' Further: ``the more the realization
of the will removes itself from the instinctual, drive-like, and
egoistic, the more it becomes a realization of the universally rational,
of the essence's law, an essence-willing, the determinateness of
objectivity, of the universally valid, in itself.'' Thus in every turn of
phrase, through every repetition, it is laid down---by a petitio
principii---that the will's universal is and shall and must be the
universal of knowledge, ``the determinateness of objectivity, of the
universally valid, in itself''; on such a petitio principii all
scientific negotiations are wasted.

That an enterprising petitio principii must be able to beget a rich
multiplicity of tautological circular proofs is understandable. By means
of the universal in general it is proved provisionally that ``the sharp
opposition'' between theoretical and practical cognition must be
abolished, in order that, once this opposition is abolished, it may be
properly proved that the universal, without difference between subjective
and objective, is the universal, the universally rational, the
objectively universal.

There is in Prof.\ Brøchner's treatment of the utterances he criticizes a
clumsiness which can easily lead a less attentive reader away from the
main point about which the inquiry turns. Concerning the marks by which I
have sought to characterize the opposition between practical and
theoretical cognition, the following is stated (p.\ 137). ``The practical
consciousness is to be personal, interested, in conflict with existence,
moved by fear and hope, full of
% --- p. 23 ---
\opage{23}\setcounter{footnote}{0}%
self-regard; the theoretical impersonal, contemplative, disinterested,
lost in the universal, moved by necessity, forgetting itself in the
grounding of the objective.'' Having enumerated these characteristic
marks, Prof.\ Brøchner then adds without further ado: ``The relative
truth in these oppositions is acknowledged through what is developed
above, but the abstract and absolute opposition is denied.'' That the
opposition I laid down between will and knowledge is denied in what the
critique has ``developed above,'' flatly denied---by a confusion of the
general with the general in general---I concede; but that anything
whatever has come forth ``through what is developed above'' that could
indicate even so much as an attempt either to efface or at least to blur
the criteria cited, I cannot possibly concede. So far is the author from
having entered upon any particular testing of the particular marks given,
that a little further on he even seems to have quite forgotten that I
have set up such marks at all. He says (p.\ 146): ``N. lays an emphasis
on this determination \emph{practical thinking}, but it nowhere in him
acquires an actual conceptual explanation, just as little as the
corresponding concept \emph{practical cognition}, and it must on the
whole be striking in what way he, during the development of his theory,
applies his psychological categories, without sharp conceptual
determination and without distinction of the essentially different.''

By seeing how matters stand with this, we obtain a not inconsiderable
contribution to a scientific appraisal of the Professor's philosophical
critique.

As proof, and that a decisive proof, that there is in actuality no
difference between practical and theoretical thinking, the following
% --- p. 24 ---
\opage{24}\setcounter{footnote}{0}%
reasoning is put forward (p.\ 146). ``\emph{Practical thinking} would most
nearly mean a reflection that determines the connection between principle
or presuppositions and consequences, between means and end, on the
practical's domain, hence, by presupposition, a formal developing of
consequences of that which is not dissolved in reflection, and this
reflection, which would yet in general acquire no place on the
paradoxically ethical's domain, would, inasmuch as the action enters,
have to be thought as absorbed, made invisible in the action's
singularity. Such a practical thinking would, however, in its \emph{form}
not distinguish itself from the theoretical thinking that draws
consequences from a not yet by cognition assimilated, e.g.\ from a merely
empirically determined, or that forms concepts of construction.'' But
what does this prove? He who concludes: All men are mortal, Caius is a
man, therefore Caius is mortal; he \emph{thinks}; he who concludes: All
men are cylinders, Caius is a man, therefore Caius is a cylinder, he
\emph{thinks}: the latter kind of thinking is
% Danish prints "er i „i sin Form“" (the "i" doubled across the line-end); translated once
``in \emph{its form}'' not to be distinguished from the former. But that
thinking in \emph{its form} is always thinking was presumably neither what
was to be proved nor what the author meant to prove.

If no weight is to be laid on the peculiarity of the principle and the
presuppositions, then there is certainly not a trace of difference
between theoretical and practical thinking: thinking in general is always
thinking. But if this is to be a proof against the essential difference
between theoretical and practical thinking, then it carries us a good
step further; we then obtain within theory at the same time a proof that
the essential difference between metaphysical and mathematical thinking
is abolished.
% Danish prints "methaphysisk" here and "methaphysiske" on p. 25 (for metaphysisk); translated as metaphysical(ly)
He who thinks metaphysically, he thinks; he who thinks mathematically, he
thinks: thinking in its
% --- p. 25 ---
\opage{25}\setcounter{footnote}{0}%
generality is always thinking; ergo mathematical thinking in \emph{its
form} is not to be distinguished from metaphysical. Will the
mathematicians hail this discovery as a new advance in science?

Metaphysical thinking is heterogeneous with mathematical as much in its
form as in its content. ``Non datur metaphysica in mathesin via,'' these
words of a keenly thinking mathematician\footnote{% Danish prints "eg" for og in the title; given correctly here
C. Ramus: Differential- og Integralregning. Preface.} have a great
testimony in the history of the sciences. ``It is only by coming forward
as mathematical science that mechanics, like the other parts of natural
science, has attained the perfection in form, the determinacy and
evidence, which is a necessary condition for every secure
predetermination of phenomena as regulated by the given forces, and so
for the possibility of successful applications. The significance of
mathematical analysis has often been misunderstood out of ignorance, both
considered as the form in which the investigation comes forward and as
the means by which it is satisfactorily carried through. It is, however,
easy to understand that analysis cannot by its own content alone bring
about the discovery of the laws and phenomena of the forces of nature. It
serves, on the other hand, to discover these phenomena's mutual
connection and their dependence on the determining
forces\footnote{C. Ramus: Analytisk Mechanik. Preface.}.''

Analytical mechanics does indeed draw consequences from something
``merely empirically determined,'' from ``a not yet by cognition
assimilated,'' insofar as it draws consequences from certain facts given
in experience; but this circumstance is very far from either lowering its
objective theoretical value or making its
% --- p. 26 ---
\opage{26}\setcounter{footnote}{0}%
thinking homogeneous with ``the reflection that determines the connection
between means and end on the practical's domain.'' The Professor's
confounding of theoretical and practical reflection everywhere betrays a
perceptible want of penetrating analysis of the matter.

Another, still stronger run at leaping the chasm between theoretical and
practical cognition Prof.\ Brøchner has already taken on p.\ 139, where
he particularly stresses ``that the will's singularity is the universal
essence's form of appearance,'' and then continues: ``In the thus
determined single the universal goes to rest only in such a way that the
single becomes transparent: the conversion is a conversion from
thinking's, and expressly from the theoretical thinking's, universality
to the personal living's singularity; but the individual as rationally
willing expresses in the singularity the universal, gives it being in its
form, which transparently reveals the universal. The difference of form
is accordingly governed by the essential unity; the content of
consciousness becomes under both forms: as content of cognition and
content of will, essentially the same. I cognize duty---and that is
again: I cognize it theoretically---e.g.\ the duty to fight for the
fatherland; I resolve to exist in the cognized, and inasmuch as the
resolve is realized, the content, as realized in the single will, is
essentially the same.'' Although the whole run simply runs out into a
confusion of the universal with the universal, even to such a degree
that, in consequence of the example set up, one is almost tempted to ask
how a man for whom the duty to fight
% Danish prints a doubtful stop after "for" ("for. Fædrelandet"); not reproduced
for the fatherland is the theory of ``the universal essence'' and ``the
will's singularity the universal essence's form of appearance'' could
have refrained from going along in the last war: I shall nevertheless, so
that it shall not seem as
% --- p. 27 ---
\opage{27}\setcounter{footnote}{0}%
if I were taking the thing all too lightly, try by way of further
elucidation to show what the result would have been if the Professor,
instead of straightway identifying the practical cognition of duty with a
theoretical cognition of the objective laws of knowledge, had taken the
marks of difference I set up under closer consideration and given himself
time to undertake an analysis of the matter.

To test the marks of difference, let one compare for once, e.g., the law
of duty with the law of gravity. On my apprehension the cognition of the
law of duty is essentially practical, the cognition of the law of gravity
on the other hand essentially theoretical; as laws they both, to be sure,
have universal validity; but the universal in the law of duty has
essentially subjective validity, the universal in the law of gravity
essentially objective validity. What this means can be seen when we
exchange the marks and make the law of gravity subjective, the law of
duty objective. The subjective moment in the cognition of the law of
gravity then lies in this, that he who is to come to clear and thorough
insight into the law's validity must \emph{be able} to think and
\emph{will} to think, and so be a thinking, willing subjectivity;
conversely, the objective moment of the cognition of duty shines forth
from this, that he who asks about duty as duty in general is not asking
about your duty or mine in particular, not about what this or that
subject happens to regard as its
% Danish prints "Piligt" (for Pligt); translated as duty
duty, but about the self-valid essence of duty. Thus we have now got
subjectivity into the cognition of the law of gravity, objectivity into
the cognition of the law of duty; but have these cognitions thereby
become homogeneous? The whole thing is empty formalism. As soon as we
stress what is peculiar in the validity of each of the two laws, the
heterogeneity must strike the eye.

The validity of the law of gravity concerns exclusively bodies as bodies,
not souls as souls, not the thinking, willing subjectivities, not the
personal in
% --- p. 28 ---
\opage{28}\setcounter{footnote}{0}%
personalities. The law that tells us that all bodies are heavy does not,
after all, tell us either that all or that only some personalities are
heavy. But when the law of gravity does not concern the essence of
personality at all, but only that of corporeality, has its fulfilment not
in the \emph{world} of personality but only in that of corporeality: then
the one who cognizes is thereby, after all, expressly called upon to look
away not merely from \emph{his} \emph{own} personality, but from \emph{all}
personality. With regard to this I call the cognition of so objective,
impersonal a law an impersonal cognition.

It is otherwise with the law of duty. The cognition may begin ever so
abstractly, ever so seemingly objectively and theoretically: the
objective and theoretical is instantly transformed into a semblance when
the emphasis falls on the validity. For which existences, then, does the
law of duty have validity? Not for the corporeal, but for the psychical,
for the thinking, willing subjectivities themselves, and so for
personalities. No single will need particularly answer for a
conscientious fulfilment of the law of gravity; it will be fulfilled all
the same. On the other hand every rational single will has particularly
to answer for a conscientious fulfilment of the law of duty; the one
single will is not, by virtue of any objective theory, to trouble itself
with answering for the other: each is in his singularity to answer
personally for himself. As a law of personality the law of duty is in
the strictest sense a law of conscience, the question of duty a question
of conscience. Is the law of conscience, then, a theory; is conscience
perhaps a science? If it stands fast that the existing personality's
knowledge with itself in conscience neither is nor can be an objective
theory; if it stands immovably fast that every question of duty
essentially is and must be a question of conscience, what instruction do
we then get out of Prof.
% --- p. 29 ---
\opage{29}\setcounter{footnote}{0}%
Brøchner's instructive words: ``I cognize duty,---and that is again: I
cognize it theoretically''?

% [p. 29 markers and the run-on sentence are in the previous fragment]

To make the distance between the practical and the theoretical in the
cognizing consciousness really vivid, let one just try to put the law of
gravity on the same footing as the law of duty and to treat the
mechanical problems as problems of conscience; the absurdity cannot then
fail to show itself. How different is not the \emph{knowledge lost in the
universal}, \emph{moved by necessity}, from the cognition that is
conditioned by \emph{self-regard} and grounded in \emph{freedom}! When Buffon
disputed with Clairaut about the universal validity of the Newtonian law,
it surely occurred to neither of the disputing parties to appeal to a
personal self-regard and say to his opponent: I feel in my conscience
that you are wrong. No, without any self-regard, without outpourings of
the heart or personal assurances, the matter was decided objectively, by
objectively necessary grounds. How strictly science deals with the
objective necessity of its grounds can be seen, among other things, from
the fact that Buffon, who held to Newton's law and to that extent was right
on the matter, was nevertheless put in the wrong in the dispute itself,
since the grounds on which he wished to rely were rejected as ``merely
metaphysical and unusable in exact science.'' How different are not the
true, irrefutable grounds of knowledge from the true, and yet so easily
contradicted, grounds of conscience! When Luther disputed with Zwingli over the
dogma of the Supper, no objective theory, be it exact or metaphysical,
was able to reconcile the disputants, and why not? Because the question
was originally not a theoretical but a practical question, not a question
of knowledge but a question of personality, in the strictest sense a
question of conscience: each of the disputing parties appealed in the last
instance, and with
% --- p. 30 ---
\opage{30}\setcounter{footnote}{0}%
full right, to his conscience. Perhaps Prof.~Brøchner will object here
that the disputants did precisely each set up his own theory, and just
thereby themselves proved that questions of conscience must in the end
seek their highest decision in theory; I need then only recall that on
neither side was it the theory that bore the personal conviction, but
conversely the personal conviction that bore the theory. Perhaps
Prof.~Brøchner will object that such decisive self-regards belong to an
antiquated standpoint, that Luther with his conscience had not yet got
beyond ``the obscure, the egoistic and drive-like,'' which ``hangs together
with the individual nature-being's unopened unity,'' to ``knowledge's
% Danish prints "uoplukkede Eenhed" (unity); Problemet S. 139 has
% "uoplukkede Enkelthed" (singularity), as quoted on p. 31. Translated as
% Nielsen prints it.
universal,'' to ``the determination of essence-universality, which
expresses itself in the being-for-itself in the form of liberated
universality that is self-consciousness's''; I need then only point out
that a ground of conscience can as little be transformed into an
objective ground of knowledge as a grounding of the law of gravity can be
transformed into a grounding of the law of duty. Should that turning point
in the history of the race ever arrive when ``the form of liberated
universality'' became one with the objective form of theory, so that one's
own personal self-regard flowed out theoretically into ``the determination
of essence-universality, which expresses itself in the being-for-itself in
the form of liberated universality that is self-consciousness's,'' then it
would have to be the turning point at which conscience itself, and with
it the law of duty, was abolished. Perhaps the Prof.\ will object that the
``objective principle of existence'' which runs through the law of duty is
of course not objective in the same way as the principle that runs through
the law of gravity, that a distinction must necessarily be made between
natural necessity and the necessity of freedom, that he has precisely
himself
% --- p. 31 ---
\opage{31}\setcounter{footnote}{0}%
admitted this essential difference by emphasizing ``the being-for-itself
in the form of liberated universality that is self-consciousness's''; I
need then only indicate the consistency with which he at one stroke knocks
his own theory to the ground. If weight is really to be laid here on the
form of liberated universality, then the difference between the necessary
and the liberated universality surely becomes a thoroughgoing difference
of principle, hence a difference of sphere, and then---the dispute is
ended.

When Prof.~Brøchner would combat ``the abstract opposition between the
universal and the particular,'' insofar as it is to condition ``the
disjunction between knowledge and existence'' (p.\ 132), he is in this
wrong only insofar as he would combat my ``theory,'' but on the other hand
entirely right insofar as he would combat his own. It is what he himself
(p.\ 139) calls ``the \emph{determination of singularity} in the will'' on
which it essentially depends. ``Is it overlooked,'' the Prof.\ continues,
``in the apprehension of this,'' namely of the \emph{determination of
singularity}, ``that the will is only will as the self-conscious, i.e.\
according to its nature universal and general essence's will, then the
will's singularity and the moment of immediacy in it is traced back to
the immediate nature-determinateness's utterance of activity in the drive,
the sensible egoistic craving, and inasmuch as this is most nearly seen
from the side through which it hangs together with the individual
nature-being's unopened singularity, which stands in immediate opposition
to knowledge's universal, then the activity that proceeds from it is
apprehended as everywhere determined by this form of singularity, and,
where it does not in general withdraw itself from consciousness, as that
which at most is reproduced by consciousness as the in its immediacy with
cognition's universal incommensurable. Is, on the contrary, this
apprehension of singularity abandoned---which would little
% --- p. 32 ---
\opage{32}\setcounter{footnote}{0}%
agree with freedom's being ascribed to the will---and is it held fast
that the will's singularity is the universal essence's form of
appearance, is \emph{spirit}-determinateness, the relation becomes
essentially another.''

Whose theory it may be that the Prof.\ is attacking here I am really
unable to guess; only so much can I say, that it is not mine. Yet it could
never have occurred to me that anyone who could reasonably be presumed
able to read would be capable of extracting so remarkable a
misunderstanding from my writings that I should actually be taken to wish
to identify the \emph{singularity} of the spirit-will with ``the immediate
nature-determinateness's utterance of activity in the drive, the sensible
egoistic craving,'' etc. How far such an apprehension departs, not only
from my way of seeing things as a whole, but expressly from my most
definite statements, how far it is from me to wish to make the existence
of the will an animal individual existence, I can show merely by citing a
couple of examples. ``No animal can thus ruin its life and miss its goal as
man can; for the animal's life is a gift of nature, man's an essential
task. As an essential task, the goal of life is an end of
self-consciousness, an ideal of freedom. Over against the ideal stands
actuality; but actuality, the naturalness that is given, must be worked
upon, altered, transformed: here there must not only be labour, here there
must be thought, here there must be will, if an ideal is to be actualized.
As the existing subjectivity that self-actively strives toward the ideal,
man is essentially spirit.'' (Forelæsninger ov. phil. Propæd. 1860---61
% Danish prints the title with a dittography across the line break,
% "Forelæs-/læsninger"; the English gives the title as intended.
p.\ 400). That the relation between ideal and actuality must originally be
clarified in cognition, I have then not overlooked either; for the
interpretation that the Critique (p.\ 144) would read into the fact that
the will can determine itself indifferently
% --- p. 33 ---
\opage{33}\setcounter{footnote}{0}%
against or even in spite of theory, ``that it should demonstrate the
inaccessibility of the will to cognition,'' does not concern me. I have
insisted as emphatically as possible on the necessity of cognition, the
practical cognition of course, of thinking, the practical thinking of
course, for the will as spirit-will. ``In
relation to actuality and the tasks of actual life, practical cognition
must absolutely have greater originality than theoretical. The will must
indeed, in order to be will, originally presuppose thought; but, for one
thing, it is, be it noted, not theoretical but practical thinking that the
will presupposes; for another, the will with this presupposing of its own
is just as original as thought. Just as there is in thought something
that cannot be derived from the will, since thought as thought is not
will, so there is also in the will something that cannot be derived from
thought, since the will as will is not thought. Certainly such categories
as freedom of choice, choice, resolution, etc.\ fall under the
illumination of thought insofar as they become clear in ethical cognition,
but in principle they are no less categories of the will; their cognition
is essentially practical.'' (Om theor. og prakt. Erkjendelse p.\ 7).

The difference between theoretical and practical cognition can be seen
expressly from the difference between conclusion and resolution.
Theoretical cognition takes shape in \emph{conclusions}, but not in
\emph{resolutions}: thought concludes, but the will resolves. The objective
conclusions of knowledge form themselves without self-regard, in
contemplative disinterestedness, without fear and hope. And why? Because
they form themselves without freedom of choice, without obligation,
without personal responsibility. In relation to the law of gravity the
knower has no freedom of choice, for the objective grounds admit of no
choice; no obligation, for
% --- p. 34 ---
\opage{34}\setcounter{footnote}{0}%
the fulfilment takes place with necessity; no responsibility, for the
question of its use does not concern knowledge: knowledge rests in
insight. It is otherwise with practical cognition; its problems are
problems of the will, its grounds grounds of the will, its conclusions
premises for a final resolution. The theoretical cognition of duty by no
means leads either to an actual fulfilment of duty or even to insight into
what fulfilment of duty means; on the contrary, it leads away from both.
``The theory'' says that the individual will is the universal, that duty
is likewise the universal, and that the will, in resolving to fulfil
duty, shows itself at one and the same time as both the particular and the
universal. These are empty abstractions. ``The theory,'' however much is
said about the concrete, never gets beyond abstraction, the abstraction of
putting the concept of the will in place of the will, the concept of duty
in place of duty, the concept of resolution in place of the actual
resolution. The actual resolution is an inward action; but precisely in
the action the will concentrates itself decisively into an individual will
with all the peculiarities and limitations of an existing individuality.
The example chosen by Prof.~Brøchner, going into battle for the
fatherland, may illuminate the matter. If for any given individual will
there were no other choice than either to go into battle or to break with
the law of duty, then there could at a pinch---although the cognition of
the law of duty, even in its pure universality, is and remains
subjective---be talk of the theoretical transparency of the resolution;
but that the duty to fight for the fatherland is a conditioned duty, a
duty dependent in its actual conditionedness on particular circumstances
and individual terms, is surely seen, among other things, from the fact
that it is not without further ado a duty for all, not e.g.\ for old men,
cripples, women and children. That the
% --- p. 35 ---
\opage{35}\setcounter{footnote}{0}%
actual conditionedness can in no possible way be delimited or determined
by theoretical cognition is evident from the fact that an individual who,
to judge by externals, is fit for battle, and to that extent obliged to
battle, can have his good and sufficient reasons why he resolves precisely
not to go into battle. The personal reasons may even be of such a nature
that, although in individual self-understanding he is fully conscious of
their validity, he nonetheless feels himself obliged not to communicate
them to anyone, but resolutely to take the responsibility upon himself and
submit to being exposed to misjudgment. Indeed, it is not even necessarily
the case that the misjudged man keeps silent about his reasons because he
carries some interesting novelistic secret or other; the ethical
explanation may simply be the consciousness that his personal reasons
would, either for ever or at least for a time, be misunderstood by others.
From another point of view something similar shows itself. Suppose that a
combatant has, in the heat of battle, been placed at a high post and
entrusted with an important charge. The task is then to do, not duty in
theoretical universality, for that no one can do---unless it were the pure
metaphysics-man---but \emph{his} duty, hence the single, determinate duty.
The task is: to obey his orders and at a decisive moment to use his own
discretion. Now, as is well known, the situation may easily take such a
shape that, if he is really to act on his own discretion, he cannot obey
his orders, or, if he is to obey his orders to the letter, he must stand
still and watch the battle being lost. Then it is in earnest a matter of
taking a resolution: here there is something to fear, and here there is
something to hope; here it shall show itself whether the individual
will---not the metaphysics-man's metaphysical one, but the actual human
being's existing individual will---has character enough to take a
responsibility upon itself.

% --- p. 36 ---
\opage{36}\setcounter{footnote}{0}%
But character! Character is surely not a theoretical something that
becomes objectively transparent to itself in an objective theory? No,
personally, character is of so subjective, practical a nature that for the
resolution itself it can use only practical cognition. Can a resolution
taken under the impression of a great responsibility perhaps be deduced
theoretically from Prof.~Brøchner's theory: ``I cognize duty---and that is
again: I cognize it theoretically\ldots.; I resolve to exist in the
% Danish prints "jer erkjender" (for "jeg erkjender", as the same quotation
% reads on p. 29); translated "I cognize".
cognized, and inasmuch as the resolve is realized, the content, as realized
in the single will, is essentially the same''? When it is a matter of
actual resolution and actual responsibility, it seems after all not so
very unreasonable to draw attention to the opposition between theoretical
and practical cognition.

It is certainly not Prof.~Brøchner's intention to abolish all difference
between cognition and will; on the contrary! he even thinks he can make me
a concession as regards the opposition. But what he holds fast most
decidedly is: that this opposition does not concern the content, but only
the form. ``Both forms: cognition and will, then become primitive
expressions for the energy of the essence revealed in self-consciousness as
unity, for the self's self-determination through which it, determining
itself, realizes itself: in both functions selfhood and energy assert
themselves under distinctive fundamental forms'' (p.\ 135). How matters
stand more precisely with these distinctive fundamental forms and the
functions corresponding to them, and how far the individual will will be
able by these functions to bring it to an actual resolution, may be seen
from statements such as these: ``The cognition that has reached its goal
would, as expression for the in-itself
% --- p. 37 ---
\opage{37}\setcounter{footnote}{0}%
transparent essence's energy, in accordance with the one essence's
necessary being in both forms, immediately convert itself into the will's
energy determined by the essence-cognition, the essence's
self-revealedness, and the will that with the most intensive energy, in
the most embracing meaning, was essence-willing: realization of the essence
% "Væsensvillen" is the form Brøchner prints (a real form, not a defect):
% essence-willing, as in the English Problemet.
according to its concretion, would be a will that immediately converted
itself into consciousness, into cognition, and into a cognition that was
determined by knowledge's universal objectively valid laws.''

If there lies hidden in these and similar statements something that really
deserves recognition, it must be something ``fine and full of spirit''
that I with my coarser sense am unable to discern; for, as far as I can
understand, there are contained herein several elements that are even very
hard to digest. In the first place, a psychological separation of two
forms distinctive of the selfhood-energy, which are not to touch the
content, and that although they are expressly opposed to each other as
primitive forms, as ``fundamental forms,'' is equally illogical and
unpsychological. A form that can be altered without the alteration
essentially concerning the content is an inessential form, accidental to
the content itself, and can by no means be a fundamental form. For what is
a fundamental form simply according to its concept? Surely a form filled
with essential content, a form whose primitive distinctiveness cannot be
altered without the content being altered. So much for the logical; now
for the psychological. The duty, e.g., to fight for the fatherland can, in
being cognized, quite certainly be apprehended and held fast by itself in
the form of ``universality and ideality,'' and thereupon in the fulfilment
be converted into the will's, i.e.\ into the ``form of singularity and
reality.'' But can this conversion now be said to be an alteration of form
that does not alter the content?
% --- p. 38 ---
\opage{38}\setcounter{footnote}{0}%
Let one just consider the \emph{same determinate} duty \emph{before} and
\emph{after} its fulfilment. Before the fulfilment the action is only a
possibility, its determinations determinations of possibility and to that
extent universal determinations. But the action is surely the content of
the duty; the content of the duty is therefore, before the fulfilment, a
content in possibility, a content of possibility. But after the
fulfilment, how everything is altered! The brave man has fulfilled his
duty, he has risked his life. When he, firmly resolved, risked his life,
he did not let his existing individual will float away, as one of the many
possibilities, in the general stream of possibilities; for he who risks in
that way is not a hero but a drunken man who is swept along in the
tumult. To risk is, ethically seen, to resolve; to resolve is to risk. In
positing the dutiful action as the decisive thing, the resolved man does
not raise himself with stoic indifference above the possible
outcome---stoic indifference is a psychological weakness that distorts the
ethical---but he lets what befalls come upon him, so that the
\emph{possible} content of the fulfilment of duty may, through the action,
the single actual action, become an \emph{actual} content. In the actual
action the will has risked itself as existing individual will, and by
putting \emph{its existence} into the content has essentially altered the
content. That Prof.~Brøchner could overlook so conspicuous an alteration
of content is strange.

Stranger still is the effort the Professor makes to explain ``the energy
of the essence revealed in self-consciousness as unity,'' ``the self's
self-determination through which it, determining itself,
% Danish prints "kvilket" (for "hvilket"; p. 36 quotes it as "hvilken");
% translated "which".
realizes itself'' in the two fundamental forms: ``cognition and will.''
The explanation is really curious. In order to assert the ``principle'' of
cognition---of objective knowledge---over against the will and to see to
it that ``the subjective essence''
% --- p. 39 ---
\opage{39}\setcounter{footnote}{0}%
may come to be active ``\emph{according to its wholeness}, its undivided
personality,'' the acute psychologist establishes no fewer than three
selfhood-energies, independent in principle, in one and the same personal
being. The first of these energies is known essentially by its having
content in itself, but no form; the two others are distinguished by each
of them being in itself a form, distinctive ``fundamental forms'' without
content. That ``the energy of the essence revealed in self-consciousness
as unity'' is in itself formless is expressly seen from the fact that, to
obtain form, it must turn first to cognition, then to the will; that
cognition and will, each by itself, must be apprehended as pure forms
without content is expressly seen from the fact that these forms are in
no way able to alter the content.

If we may now take this acute division in earnest, literally, according to
the determinations of concept stated, then we do indeed arrive at a
dualism that is quite otherwise harsh than the one I pointed out. Is
cognition really independent of the will in the sense that it gets its
selfhood-energy, not from the will as will, but from the first
self-determination, independent of the will, that as yet formless,
mystical potency which originally seeks its form in cognition: then we
surely have a whole region of knowledge, a circle of objective theories, of
which it can truly be said that it stands independent of the will, wholly
independent with respect to content as well as to form. For objective
knowledge can then by no means depend on the will as regards its content,
when the content comes from the formless potency; nor can knowledge depend
on the will as regards its form, when the form of cognition has not merely
as great a primitivity as the form of the will, but an even greater one.
% --- p. 40 ---
\opage{40}\setcounter{footnote}{0}%
If this is not a psychological dualism, what then is dualism?

Perhaps we may not lay such weight on the separation of form and content.
We will then try it another way and assume that the whole distinguishing
between three different kinds of selfhood-energy does not mean much; for
by assuming this we seem, at least in some places, to come nearer to the
author's meaning. We are told that ``the cognition that has reached its
goal'' would immediately convert itself into the will's energy determined
by the essence-cognition; we are told that ``the willing that with the
most intensive energy was essence-willing'' would be ``the will that
immediately converted itself into consciousness, into cognition,'' and
that into the bargain into a cognition that was determined by knowledge's
universal, objectively valid laws. From what is thus said here we get the
notion that it is ``mediation,'' the old well-known speculative mediation,
that is to rise again and keep step with the reasoning critique. But the
critique and mediation cannot be reconciled. If mediation is to rule, it
will in no way do to assume an essential alteration of form without a
corresponding alteration of content. Say what one will about mediation: in
the chapter on form and content it allows no halving; its thinking is
immanent, and immanent thinking does not engage in any critical
parcelling-out of form and content. How then should we be able to conceive
of the cognition that has reached its goal going about converting
itself---immediately---into will? The goal of cognition is surely nothing
but objective theory, nothing but science. As long as objective cognition
has not exhausted all possible material of knowledge, it must of course go
forward from knowledge to knowledge without converting itself into will.
Since, now, the author will not admit any
% --- p. 41 ---
\opage{41}\setcounter{footnote}{0}%
absolute limit of knowledge, cognition will of course never reach its
goal, and the expected conversion will never occur.
% [p. 41 markers and the run-on sentence are in the previous fragment]
Suppose that cognition, after having reached a relative goal, converts
itself into will: what then halts cognition; what is it that brings about
the conversion? According to mediation, cognition would have to halt of
itself, of itself pass over into will, because in its essence it is will.
But that the critique forbids; the reasoning critique has expressly posited
the form of cognition by itself, independent of the form of will. According
to sound common sense it would have to be the will that halted cognition and
brought about the conversion; but that conflicts with the theory put
forward, which posits cognition as independent of the will. It must, then,
be that selfhood-energy formless in itself, that mystical potency which is
neither cognition nor will, that has to convert cognition into will. But
even when this is assumed, there is still no reason here why cognition, ``as
expression for the in-itself transparent essence's energy, in accordance
with the one essence's necessary being in both forms,'' should ``immediately
convert itself'' into the form of will. The only reason the explanation
contains is a sort of reason why ``the in-itself transparent essence's
energy,'' after having expressed its content in the form of cognition, can
change over and express its content in the form of will. But such a change of
form is then very far indeed from signifying the transition of the one form
into the other.

If we try to think the transition by starting from the will's
self-conversion into knowledge, we are no nearer. What sort of will is it
that --- with ``the most intensive energy'' --- immediately converts itself
into consciousness? According to the author's own assumption, the will as
will cannot, after all, be without consciousness; the will with the ``most
intensive energy'' must therefore be a will with the
% --- p. 42 ---
\opage{42}\setcounter{footnote}{0}%
most intensive consciousness. But when the will with the most intensive
energy itself, namely in the will's fundamental form, \emph{is} the most
intensive consciousness, how then can it \emph{convert} itself into
consciousness, into what it already is? Here there is, after all, no
conversion. Is the wonderful conversion perhaps to consist in this, that the
will, after having for a time kept itself, with a less intensive energy, in
the form of ordinary will, simple, practical will of character, at last
reaches the most intensive energy, and then suddenly, by means of the most
intensive energy, veers round --- ``immediately,'' as when it is said that
the wind veers round --- and transforms itself from will into not-will, into
cognition, into objective cognition, cognition of knowledge's universal,
objectively valid laws: then the whole explanation will come very close to
the meaningless. The philosopher who could make such a conversion clear to
us in a definite psychological example, such an immediate objective
veering-round in the subjective, would have to be a true theoretical
sorcerer, an objective psychological conjurer.

One last way out still remains; that too must not be left untried. In order
to maintain the undivided unity of personality with itself in cognition and
will, the Professor appeals to what he calls ``the energy of the essence
revealed in self-consciousness as unity,'' to ``the self's
self-determination through which it, determining itself, realizes itself,''
and so on. It is therefore the relation between the will and the self's
self-determination on which attention must fix itself. According to the
customary way of seeing things one might be tempted to assume that the
self's self-determination was properly its original willing, and the will
itself the centre of personality; but this is by no means the case. ``The
energy of the essence revealed in self-consciousness as unity''
% --- p. 43 ---
\opage{43}\setcounter{footnote}{0}%
has, when we abstract from the pseudo-independent self-determination of the
energy of cognition --- mirabile dictu --- two kinds of specially
determining self-determination ``through which it, determining itself,
realizes itself.'' The first self-determination is the one formless in
itself --- the mystical potency; the latter, the one empty of content in
itself, is the will. The will as will is therefore not the self-determining
unity of personality; for the will first comes forth when the first,
formless self-determination determines itself to lay the content of
consciousness into the other determining self-determination, namely the one
that determines the form. The will, seen in this light, is then as good as
not accountable. The will, alas, the poor will, has been taken captive by a
mystical potency and cannot originally determine itself: its
self-determination is each time dependent on whether the self's first
self-determination, ``through which it, determining itself, realizes
itself,'' will determine it to self-determination.

It is on such a foundation that Professor Brøchner would maintain the
undivided unity of personality! First he makes, on his own theoretical
plenary authority, the existing consciousness into a scientific theory; next
he confounds the subjective validity of the universal with the objective,
even to such a degree that he does not see himself able to make a
distinction between grounds of science and grounds of conscience, between a
cognition of the law of duty and a cognition of the law of gravity.
Next, after the spirit-man existing in flesh and blood has been confused
with the silhouette of a metaphysics-man bleached out by abstractions, the
universal concept of the individual will is by a mistake put in the place of
the will struggling for its independence. Finally, after the subjective
universal essence with the three selfhood-energies, the one with content
without form, the other two
% --- p. 44 ---
\opage{44}\setcounter{footnote}{0}%
with form without content, has become so objectively transparent to itself
that subjectivity's original self-determination can in perfect transparency
determine itself to the will's self-determination, the goal is reached:
action is theory and theory action, the difference at all events only a
difference of form that does not alter the content. On such a foundation it
is now an easy matter to deny the absolute boundary of knowledge and reject
the mystery; on such a foundation a criticizing philosopher can then easily
prove both the impossibility of God's personality and the impossibility of
the validity of prayer and the impossibility of the soul's immortality.

On such a foundation, an airy substratum of essence-empty abstractions,
skewed deductions, smeared-out turns of phrase and short-winded remarks,
Professor Brøchner has now really stilted himself up into a philosophical
authority, a high priest of knowledge, who purifies faith, preserves ``the
religious'' by ``abolishing the religions,'' and proclaims the gospel of the
metaphysics-man, of the ``universal human essence,'' in that knowledge is
chosen. Thus run his words: ``For inasmuch as knowledge is chosen, there is
chosen that which, as essence-expression for the universal human essence,
bears in itself the possibility of concluding itself into a totality with
inner infinity; that which, inasmuch as it has the ideal primacy in the
indivisible human spirit, can become the norm for an ethical life-whole;
that which, while it sees through the finite religious forms' essence,
comprehends them in their origin out of the self-unfolding, toward its
principle striving human nature, and therefore does not retain in them a
disturbing barrier of consciousness, and which, inasmuch as it in the
historical religious forms separates the transient and finite from the
abiding, essential, and eternal, the non-religious from the religious, can
stand in harmony with and take up into itself
% --- p. 45 ---
\opage{45}\setcounter{footnote}{0}%
this essential, the religious according to its true and eternal
significance.'' (pp.\ 225--26).

If someone now, after having laboriously worked his way through the
Professor's book \emph{Problemet om Tro og Viden}, were to declare outright
that the reason why Prof.\ Brøchner has on this occasion come forward
against the last possible attempt to defend the faith of revelation against
the attacks of science is plainly not love of science, but a smouldering
hatred of positive Christianity, would one not then answer back vigorously
and raise a hue and cry: ``No insinuations, no personalities! Test the
proofs, but do not mutter about the motives!'' I assume that one would do
so, and rightly. But what does that prove? It proves, after all, precisely
what I maintain, but what Prof.\ Brøchner will not concede: that there is a
difference, a decided difference of principle, between questions of
knowledge and questions of conscience, between grounds of science and
grounds of personality. It proves precisely what I maintain, but what Prof.\
Brøchner expressly denies: that there is a difference, an incontrovertible
difference of principle, between the clear insight into an objective theory
% Danish prints "imellen" (for imellem); translated as between
and the ``mirroring'' of that interior of personality which is never
exhausted in any scientific concept.

Two knowers can become objectively transparent to each other in the same
knowledge of the same science; but two personalities can never become
objectively transparent to each other in the same personal concerns; for
this stands fast: the deeper the personality, the richer the inwardness; but
no science owns the magic key that opens the door of inwardness.

But although there are depths in existing personalities into which the
objective light of knowledge never reaches,
% --- p. 46 ---
\opage{46}\setcounter{footnote}{0}%
it by no means follows from this that a general, approximately grounded
judgment could not be passed on personal actions, on actual characters and
their greater or lesser agreement as well with themselves as with the
universally human; only that the apprehension, as has been said, becomes a
% Danish prints "phychologisk" (for psychologisk); translated as psychological
psychological mirroring, not an exhaustive comprehension. With regard to
this, then, I can very well, without coming too near the personal, point out
that the interest which comes to the fore in Prof.\ Brøchner's critique
before us is not a positive interest of knowledge, but a negative interest
of faith. I discern this from the following indications. If it had really
been Prof.\ Brøchner's intention to assert the right of science against
encroachment from the standpoint of faith, then he would, on this point at
least, have had to side with me. The discussion between us would then have
become a scientific discussion about something purely scientific, e.g.\
about the problems in Grundideernes Logik. Of this work there lies before
the public a series of analyses running to no fewer than 60 sheets. Here,
then, whatever the result might be, there would have been something to
investigate. That, alongside these purely scientific endeavours, I thought I
could moreover preserve and defend the old faith of revelation, the
Professor could, after all, for the time being have let stand as, in his
eyes, a psychological curiosum, a ``mirroring'' of something personal that
did not quite let itself be exhausted in the objective concept, at least
not until I had got so far that I could present my philosophy of religion as
a worked-out whole. But now Prof.\ Brøchner has gone precisely the opposite
way. Far from waiting for an authentic and coherent presentation of my
doctrine of faith, he has in great haste thrown himself upon ``remarks,''
``reports,'' and ``aphoristic
% --- p. 47 ---
\opage{47}\setcounter{footnote}{0}%
utterances,'' has treated every sentence, every word, every jot, urging,
specifying, distinguishing, reducing, conceding, restricting, reasoning, and
so on, with a zeal as systematic as if he had had a finished system before
him. This haste, this zeal plainly points to the chief task being to do
away with ``the faith of revelation'' and get ``the religions'' abolished in
order to be able to preserve ``the religious.''

That the interest of knowledge does not stand in the foreground I can
easily prove. I need, after all, only refer the reader to the assessment of
Grundideernes Logik found in the Professor's book (pp.\ 154--160). For a
grounded judgment on this assessment it is not at all necessary that the
educated reader be particularly at home in metaphysics and know his way
either about the relation between system and schema or about ``the power''
or about ``the antilogical''; if he only has a tolerably clear notion of what
an author who has at least set out to achieve something scientific is
entitled to demand of the one who would submit his work to a scientific
examination, he will without difficulty grasp what everyone who has respect
for science must judge of Prof.\ Brøchner's scientific assessment. For it is
certain that a half-metaphysical female, if she only knew how to use such
words as ``subreption,'' ``surreptitious insertion,'' ``unclarity,''
% Danish prints the „ before "Vilkaarlighed" malformed (",."); English normal
``arbitrariness,'' ``sophistry,'' ``metaphor,'' ``category,'' ``not
deduced,'' and the like, in places where they ring well, could in a
newspaper article deliver an equally thorough, equally scientific
assessment, and on top of that make it quite readable.

That Grundideernes Logik is not to Professor Brøchner's taste must no doubt
be explained by several reasons, among others by this, that it does not at
all suit his critical
% --- p. 48 ---
\opage{48}\setcounter{footnote}{0}%
talent. ``It is an old objection that philosophy hangs on words, that its
explanations are verbal explanations, its disputes disputes about words, its
wisdom verbal junk. This objection is not wholly without ground when one
one-sidedly fixes attention on the shadow sides, on the barren subtleties,
the tasteless excesses in empty verbal cleverness, over which a John of
Salisbury, an Erasmus of Rotterdam, and with us a Holberg have, each in his
age, held court against the clients of speculative knowledge.'' With these
words begins the second part of the logic in question, and perhaps there is
something repellent in them. Certainly weight is laid on the word.
``Philosophy is the science of the word; it must necessarily hold on to the
word, as surely as it will hold on to the thought and hold fast the concept.
Just as a soul cannot be without a body, so neither can the concept, as the
concept determinate in itself, be without a determining word; for if the
thought were not, then there were no word, and if the word were not, then
there were no thought either.'' But in the further development it is
nevertheless stressed with decisive emphasis that the development of thought
conditioned by the word ought always to go hand in hand with a
corresponding development of the case. Here now, if I am not mistaken, is
the real point of collision between Grundideernes Logik and Prof.\
Brøchner's talent.

That Prof.\ Brøchner's talent is a critical talent of a peculiar kind,
namely a talent for philosophical criticism of words, shall not in this
connection be tossed off as a passing assertion, but, by way of conclusion,
be made evident by a couple of striking examples. In order to illuminate
what he calls ``the theory's consequences with respect to the
psychological'' (p.\ 178 ff.), the Professor begins with a couple of lengthy
citations taken from my book ``Om den gode Villie,'' and then introduces the
illumination itself with
% --- p. 49 ---
\opage{49}\setcounter{footnote}{0}%
the following declaration: ``We have reproduced these developments, also
the `lazzi' placed toward the end, so completely, because only such
completeness can give a clear picture of the way in which the theory must be
defended'' (p.\ 183). By this forced run-up the reader is naturally to be
impressed into assuming that what I have said is not only some insipid and
meaningless twaddle, but on top of that something most unseemly. My words
are these: ``Does it then stand fast that faith can in no way be
transformed into knowledge or knowledge into faith, then it also stands fast
that faith can just as impossibly be halved by knowledge as knowledge by
faith. Can a consciousness not be halved by friendship and mathematics, then
it can still less be halved by ethics and astronomy. There are mathematical
friends, but there is no friendly mathematics; there are enamored chemists,
but no chemical being-in-love, no enamored chemistry. There can be ethical
and unethical astronomers, but astronomy is neither ethical nor unethical.
Science is in other words neither ethical nor unethical, neither believing
nor unbelieving; love is practical, but science is theoretical.'' It is only
a few lines that are in question here, and yet I concede that he who can
shake the fundamental thought compressed in these few lines, where a
development of the case is concerned, overturns my whole view. ``But'' ---
someone perhaps objects --- ``Prof.\ Brøchner has indeed shaken it!'' I
% Danish prints a stray ;-shaped mark before „Hr.; not rendered
answer: Professor Brøchner does not, as any tolerably attentive reader can
see, go into the case at all; no, he fiddles with the words: his
philosophical critique is in this as in other very important respects a
criticism of words, and a rather pusillanimous one at that. One reads
(p.\ 187): ``The `lively' analogies N. adds for elucidation are most nearly
only a not very successful attempt at the burlesque, and prove,
% --- p. 50 ---
\opage{50}\setcounter{footnote}{0}%
with respect to the matter treated, if possible, less than
nothing.''\footnote{Had Prof.\ Brøchner, with regard to ``the mathematical
friends'' and ``the friendly mathematics,'' really elucidated, or even
merely made an attempt to elucidate, something that concerned ``the matter
treated,'' then we might perhaps have received an interesting elucidation of
how the personal relation between mathematical friends can contribute
\emph{either to weakening or to strengthening the objective mathematical
% Danish prints "mathemathiske" (for mathematiske); translated as mathematical
proofs}, how, in other words, grounds of personality can be transformed into
grounds of science, an elucidation that would surely not fail to bring about
a great revolution in mathematics as well as in the other sciences. But
instead of elucidating something about the matter, the Professor has, perhaps
on personal grounds, preferred to elucidate something about the words. The
elucidation is ``that one cannot without further ado, because an adjective
can be predicated of a substantive, therefore, conversely, predicate an
adjective corresponding to the substantive of a substantive corresponding to
the adjective,'' an elucidation which, as far as the matter goes, stands about
on a level with Per Degn's elucidation of Imprimatur: ``All the words that
can be named are referred to 8 things, which are Nomen, Pronomen, Verbum,
Participium, Conjugatio, Declinatio, Interjectio.''

To what degree Professor Brøchner can forget the case merely in order to
criticize the words, and how strict he is when he criticizes the words, is
seen from this, that his acuteness has been able to discover ``the
burlesque'' even in the examples cited. Propriety is a virtue, even in a
philosophical exposition. But since the boundary of the proper is, after
all, not objectively scientific, the philosophizing personalities can
nonetheless also push the demand for propriety so far that the exposition
suffers for it. If the Professor sees himself able to prove in ``a
theoretically exhaustive determination of the concept'' that I have in the
examples cited mixed in ``impure ingredients'' which can be brought under
the category of ``the burlesque,'' then I must of course bow to objective,
incontrovertible grounds; but until such a proof is presented, I reserve my
personal right. For it is really my conviction that philosophy, even if it
% note continues on p. 51
were personified in the pure metaphysics-man and by this transformation ---
which of course would not be permitted to touch the content, but only the
form --- had become as proper as an old schoolmistress: the old maid, unless
she had a hysterical fit of crossness, would still have trouble finding
``the burlesque'' in my innocent examples.}

% --- p. 51 ---
\opage{51}\setcounter{footnote}{0}%
Even in the little piece (pp.\ 184--186) where the criticizing author, after
sundry circumlocutory remarks, at last acts as though he were going into the
case, the critique nevertheless continues to be a criticism of words, in
which one skewed remark relieves the other. ``Inasmuch as knowledge
acknowledges its boundary, it acknowledges in that which is faith's sphere a
mystery for knowledge. Inasmuch as faith has a boundary against knowledge,
that which is knowledge's object becomes a mystery for faith. Faith is
accordingly knowledge's mystery, knowledge faith's, and the mystery, which,
thus apprehended, has nothing mysterious, in this way does not prevent the
restriction, but precisely expresses it.'' This example of an empty play on
words may figure here as instar omnium; only one must grasp ``the fine
point.'' ``The fine point'' is that a mystery, when one knows what it is: ``a
boundary against knowledge,'' is no longer a mystery; for one knows, after
all, what it is: ``a boundary against knowledge.'' So a secret, when I know
what it is, namely that which I cannot know, is no longer any secret, for I
know, after all, what a secret is, namely: that which I cannot know.

The following remark is perhaps not quite so empty, but it is both crooked
and skewed. ``Knowledge is, by understanding the \emph{intermediate
determinations hidden in the mystery}, to be able to think the
power-revelation, `in which the believer sees the miracle,' as subsisting
with the laws' immutability, but this it can only
% --- p. 52 ---
\opage{52}\setcounter{footnote}{0}%
when it, while the miracle's name is preserved, lets the miracle as
\emph{miracle vanish in a supposed rational}.'' Readers who wish to know my
view of the miracle in a corresponding development of the case I refer to my
Forelæsninger om Hindringer og Betingelser for det aandelige Liv i Nutiden
(IX and X); about the remark in the criticism of words before us I can be
brief. What is characteristic of the miracle is not at all that a breach is
made in the laws, but that a fact is immediately referred to God's almighty
will. That the fundamental relation between the realm of nature and God's
almighty will is a relation of inner infinity, therefore a relation in which
the will works as one with the laws, can be seen in all generality, without
the fundamental relation itself in its actual articulations therefore
ceasing to be an absolute mystery for knowledge as well as for faith.
Whether such a mystery is to be accepted or denied is not a question of
science; no, it is now as before a question of conscience.

An uncritical confusion of questions of conscience with questions of
science is surely the strongest testimony against the scientific character
of Prof.\ Brøchner's philosophical critique.

\begin{center}\rule{3cm}{0.4pt}\end{center}

\end{document}
